\def\titlestr{Kepler's Ellipses}
\ifx\makebook\undefined
	% ----------------------------------------------------------------
% ----------------------------------------------------------------
% AMS-LaTeX Paper ************************************************f
% **** -----------------------------------------------------------
% documentclass[fleqn]{report}
%\RequirePackage{silence}
%\WarningFilter{lualatex}{Token not allowed}
% Einzelne Warnung eines Pakets unterdrücken:
%\WarningFilter{pdftex}{destination with}

% Alle Warnungen eines bestimmten Pakets unterdrücken:
%\WarningsOff[lualatex]
\PassOptionsToPackage{unicode}{hyperref}
%\PassOptionsToPackage{naturalnames}{hyperref}
\documentclass[a4paper,fleqn]{article}
%\documentclass[a4paper,fleqn]{scrreprt}
%\documentclass{styles/wissdoc-kw-eng}
\setlength{\oddsidemargin}{0.0cm}
\setlength{\evensidemargin}{0.0cm}
\setlength{\paperheight}{297mm}
%\usepackage[utf8]{inputenc}
\usepackage{german}
\usepackage{graphicx} 
\usepackage{amsmath}
\usepackage{amssymb}
\usepackage{alltt}
\usepackage{color}
\usepackage{fancyvrb}
\usepackage{verbatim}
\usepackage{mathtools}
\usepackage[compat=1.1.0]{tikz-feynman}
\DeclarePairedDelimiter\bra{\langle}{\rvert}
\DeclarePairedDelimiter\ket{\lvert}{\rangle}
\DeclarePairedDelimiterX\braket[2]{\langle}{\rangle}{#1\,\delimsize\vert\,\mathopen{}#2}
\DeclarePairedDelimiterX\expect[1]{\langle}{\rangle}{#1}
%\newcounter{thm}[section]
%\newtheorem{theorem}[thm]{thm}
%\newtheorem{lemma}[theorem]{Lemma}
%\newtheorem{proposition}[theorem]{Hilfssatz}
%\newtheorem{corollary}[theorem]{Corollary}
%\newtheorem{definition}[theorem]{Definition}
%\newenvironment{theorem}[Theorem]{\begin{trivlist}
%\item[\hskip \labelsep {\textit \textit{#1}}]}{\end{trivlist}}
%\newenvironment{proof}[1][Beweis]{\begin{trivlist}
%\item[\hskip \labelsep {\bfseries #1}]}{\end{trivlist}}
%\item[\hskip \labelsep {\bfseries #1}]}{\end{trivlist}}
\newcommand*\codeRed[1]{\textcolor[rgb]{1,0,0}{\textbf{#1}}}
\newcommand*\codeGreen[1]{\textcolor[rgb]{0,0.5,0}{\textbf{#1}}}
\newcommand{\lb}{[}
\newcommand{\rb}{]}
\newcommand{\codeCall}[1]{\textbf{\texttt{#1}}}
\newcommand{\mcL}{\mathcal{L}}
\newcommand{\kb}{k_{\mathrm{B}}}
%\newcommand{\titlemath}[1]{\texorpdfstring{#1}{}}
\newcommand{\titlemath}[2]{\texorpdfstring{#1}{#2}}
\newcommand{\pmat}[1]{\begin{pmatrix}#1\end{pmatrix}}


%\newenvironment{\vspace{1ex}\noindent{\bf Proof}\hspace{0.5em}}
%	{\hfill\qed\vspace{1ex}}
 

\newenvironment{example}[1][Example]{\begin{trivlist}
\item[\hskip \labelsep {\bfseries #1}]}{\end{trivlist}}
\newenvironment{remark}[1][Remark]{\begin{trivlist}
\item[\hskip \labelsep {\bfseries #1}]}{\end{trivlist}}

%\newcommand{\qed}{\nobreak \ifvmode \relax \else
%      \ifdim\lastskip<1.5em \hskip-\lastskip
%      \hskip1.5em plus0em minus0.5em \fi \nobreak
%      \vrule height0.75em width0.5em depth0.25em\fi}
%\newcommand{\qed}{$\square$}
\newcommand{\qt}{\texttt}

\newcommand{\bibdef}[4]{\bibitem{#1} {\bf #2:} {\it #3}\\#4}


\newcommand{\gauss}[1]{\lfloor #1\rfloor}
\numberwithin{equation}{section}


%\setcounter{equation}{1}

\setcounter{tocdepth}{4}
\setcounter{secnumdepth}{4}

\newcommand{\ltri}{\vartriangleleft}
\usepackage{amsthm}
\newtheorem{theorem}{Theorem}[section]
\newtheorem{definition}{Definition}[section]
\renewcommand{\proofname}{Beweis}
\definecolor{dblue}{RGB}{150,50,150}
%\usepackage[colorlinks=true, linkcolor=dblue, urlcolor=cyan]{hyperref}
\usepackage{hyperref}
%\usepackage{cleveref}

%\usepackage{crossreftools}

%\addto\captionsngerman{%
 \renewcommand{\abstractname}{Vorbemerkung}
%}
	\begin{document}
	\title{\titlestr}
	\author{Hugo Meder}
	\maketitle
	\tableofcontents
	\newpage
	%\begin{abstract}
	Dieser Text ist im Rahmen meiner Vortragsreihe 'Ästhetik der Rationalität' entstanden. Es handelt sich \textbf{nicht} um den Text des Vortrags selbst. Er enthält aber die Mathematik und die Physik die im Vortrag dargestellt wird. Ziel ist ein doppeltes:
\begin{itemize}
	\item Dem Interessierten Zugang zu einer formal exakten Formulierung des Inhalts des Vortrags zu gewähren.
	\item Für mich selbst erfüllt das Schreiben den Zweck mir die Struktur und die logischen Zusammenhänge klar zu machen. Diese Struktur führt dann auch dazu Reihenfolge und Darstellung des Vortrags zu gestalten.
\end{itemize}

	%\end{abstract}
	\newpage

\else
	\begin{chapter}{\titlestr}
\fi

\begin{section}{The First Two Kepler Laws}
The first of Kepler's laws states that planets move on elliptical orbits with the Sun located at one of the focal points.
If we place the Sun at the origin of coordinates and the major semi-axis on the $x$-axis, we obtain
\begin{eqnarray*}
z(\varphi) &=&
\begin{pmatrix}
x(\varphi)\\
y(\varphi)
\end{pmatrix}
=
\begin{pmatrix}
a\cos(\varphi)+f\\
b\sin(\varphi)
\end{pmatrix}
\end{eqnarray*}
where $\varphi$ plays the role of a curve parameter. $a$ and $b$ are the lengths of the major and minor semi-axes, f is the distance from the focal point to the center of the ellipse. It holds:
\begin{eqnarray*}
f^2 &=& a^2-b^2
\end{eqnarray*}
The second law states that the line connecting the Sun and the planet sweeps out equal areas in equal time intervals. The swept area as a function of $\varphi$ is:
\begin{eqnarray*}
A(\varphi) &=& \frac{b}{2}\left(a\varphi+f\sin(\varphi)\right)
\end{eqnarray*}
So it must hold:
\begin{eqnarray*}
\dot{A} &=& A^\prime(\varphi(t))\dot{\varphi}(t) = w
\end{eqnarray*}
for a constant $w$. In other words:
\begin{eqnarray*}
\dot{\varphi}(t) &=& \frac{w}{A^\prime(\varphi(t))} = \frac{2w}{b}\frac{1}{a+f\cos(\varphi(t))}
\end{eqnarray*}
\end{section}

\begin{section}{The Acceleration}
\label{ell_accel1}
The velocity is obtained as:
\begin{eqnarray*}
\dot{z} &=& z^\prime(\varphi(t))\dot{\varphi}(t)
\end{eqnarray*}
Accordingly, the acceleration is:
\begin{eqnarray}
\label{abstracc}
\ddot{z} &=& z^{\prime\prime}(\varphi(t))\dot{\varphi}^2(t)+z^\prime(\varphi(t))\ddot{\varphi}(t)
\end{eqnarray}
To compute $\ddot{\varphi}(t)$, we treat $\dot{\varphi}$ as a function of $\varphi$:
\begin{eqnarray*}
\dot{\varphi}(\varphi) &=& \frac{2w}{b}\frac{1}{a+f\cos(\varphi)}
\end{eqnarray*}
Then
\begin{eqnarray*}
\ddot{\varphi}(\varphi(t)) &=& \dot{\varphi}^\prime(\varphi(t))\dot{\varphi}(\varphi(t))\\
&=& \frac{2w}{b}\frac{f\sin(\varphi(t))}{(a+f\cos(\varphi(t)))^2}\frac{2w}{b}\frac{1}{a+f\cos(\varphi(t))}
\end{eqnarray*}
Summarized we then obtain:
\begin{eqnarray*}
\ddot{\varphi}(\varphi(t)) &=& \left(\frac{2w}{b}\right)^2\frac{f\sin(\varphi(t))}{(a+f\cos(\varphi(t)))^3}\\
\dot{\varphi}^2(\varphi(t)) &=& \left(\frac{2w}{b}\right)^2\frac{1}{(a+f\cos(\varphi(t)))^2}
\end{eqnarray*}
To simplify the further calculation, we introduce a common factor $\kappa(\varphi(t))$:
\begin{eqnarray*}
\kappa(\varphi(t)) &=& \left(\frac{2w}{b}\right)^2\frac{1}{(a+f\cos(\varphi(t)))^3}
\end{eqnarray*}
And then obtain:
\begin{eqnarray*}
\ddot{\varphi}(\varphi(t)) &=& f\sin(\varphi(t))\kappa(\varphi(t))\\
\dot{\varphi}^2(\varphi(t)) &=& (a+f\cos(\varphi(t)))\kappa(\varphi(t))
\end{eqnarray*}
Now expression \ref{abstracc} can be easily evaluated:
\begin{eqnarray*}
\ddot{z}(t) &=& z^{\prime\prime}(\varphi(t))\dot{\varphi}^2(t)+z^\prime(\varphi(t))\ddot{\varphi}(t)\\
&=& \left(z^{\prime\prime}(a+f\cos(\varphi(t)))+z^\prime(\varphi(t))f\sin(\varphi(t))\right)\kappa(\varphi(t))\\
&=&\left(
\begin{pmatrix}
-a\cos(\varphi(t))\\
-b\sin(\varphi(t))
\end{pmatrix}
(a+f\cos(\varphi(t)))+
\begin{pmatrix}
-a\sin(\varphi(t))\\
b\cos(\varphi(t))
\end{pmatrix}
f\sin(\varphi(t))\right)\kappa(\varphi(t))\\
&=&\begin{pmatrix}
-a\cos(\varphi(t))(a+f\cos(\varphi(t))-af\sin^2(\varphi(t))\\
-b\sin(\varphi(t))(a+f\cos(\varphi(t))+fb\cos(\varphi(t))\sin(\varphi(t))
\end{pmatrix}
\kappa(\varphi(t))\\
&=&\begin{pmatrix}
-af-a^2\cos(\varphi(t))\\
-ab\sin(\varphi(t))
\end{pmatrix}
\kappa(\varphi(t))\\
&=&-a\begin{pmatrix}
f+a\cos(\varphi(t))\\
b\sin(\varphi(t))
\end{pmatrix}
\kappa(\varphi(t))\\
&=&-a\kappa(\varphi(t))z(\varphi(t))
\end{eqnarray*}
That is, the acceleration is $-a\kappa(\varphi(t))$ times the distance vector $z(\varphi(t))$. Now it is helpful to know that
\begin{eqnarray*}
|z(\varphi)| &=& a+f\cos(\varphi) = r
\end{eqnarray*}
\begin{proof}
\begin{eqnarray*}
|z(\varphi)|^2 &=& x^2+y^2 = \left(f+a\cos(\varphi)\right)^2+\left(b\sin(\varphi)\right)^2\\
&=& f^2+2fa\cos(\varphi)+a^2\cos^2(\varphi)+b^2\sin^2(\varphi)\\
&=& f^2+2fa\cos(\varphi)+a^2\cos^2(\varphi)+b^2\left(1-\cos^2(\varphi)\right)\\
&=& f^2+b^2+2fa\cos(\varphi)+\left(a^2-b^2\right)\cos^2(\varphi)\\
&=& a^2+2fa\cos(\varphi)+f^2\cos^2(\varphi)\\
&=& \left(a+f\cos(\varphi)\right)^2\\
|z(\varphi)| &=& a+f\cos(\varphi)
\end{eqnarray*}
\qed
\end{proof}
Therefore the absolute value of $\ddot{z}$ becomes:
\begin{eqnarray}
\label{ell_acc}
|\ddot{z}(t)| &=& a\kappa(\varphi(t))|z(\varphi(t))|=a\left(\frac{2w}{b}\right)^2\frac{1}{(a+f\cos(\varphi(t)))^3}|z(\varphi(t))|=a\left(\frac{2w}{b}\right)^2\frac{1}{r^3(t)}r(t)\cr
&=& a\left(\frac{2w}{b}\right)^2\frac{1}{r^2(t)}
\end{eqnarray}
$w$ can be expressed in terms of the orbital period $T$. After time $T$ the entire area of the ellipse is swept:
\begin{eqnarray*}
wT &=& \pi ab\\
w &=& \frac{\pi ab}{T}
\end{eqnarray*}
Thus we obtain for the magnitude of the acceleration:
\begin{eqnarray}
\label{accel}
|\ddot{z}(t)| &=& \left(2\pi\right)^2\frac{a^3}{T^2}\frac{1}{r^2(t)}
\end{eqnarray}
\end{section}

\begin{section}{The Third Kepler Law}
\label{ell_accel2}
It states that for two different planets with major semi-axes $a_1$ and $a_2$ and orbital periods $T_1$ and $T_2$ it holds:
\begin{eqnarray*}
\frac{a_1^3}{T_1^2} &=& \frac{a_2^3}{T_2^2}
\end{eqnarray*}
Thus the dependence in \ref{accel} on $a$ and $T$ can be eliminated and written as
\begin{eqnarray*}
|\ddot{z}| &=& \frac{K}{r^2}
\end{eqnarray*}
\begin{eqnarray}
\label{kdef}
K&=&\left(2\pi\right)^2\frac{a^3}{T^2}
\end{eqnarray}
is a constant fixed for the solar system. We thus obtain the following law of motion:
\begin{eqnarray}
\label{kepler_motion}
\ddot{z}(t) &=& -\frac{Kz(t)}{|z(t)|^3}
\end{eqnarray}
This law of motion now holds for \textbf{all} planets of our solar system.
\end{section}
\ifx\makebook\undefined
	\end{document}
\else
	\end{chapter}
\fi